\documentclass[11pt]{article}
\usepackage[T1]{fontenc}
\usepackage[utf8]{inputenc}
\usepackage{lmodern}
\usepackage[margin=1in]{geometry}
\usepackage{amsmath,amssymb,amsthm,mathtools}
\usepackage{booktabs,array,longtable}
\usepackage{enumitem}
\usepackage{hyperref}
\newcommand{\OPHPaperReleaseID}{r1374}
\newcommand{\OPHPaperReleaseDate}{April 4, 2026}
\newcommand{\OPHPaperReleaseBanner}{%
\begin{center}
\small\textbf{Paper release:} \texttt{\OPHPaperReleaseID}\quad\textbf{Released:} \OPHPaperReleaseDate
\end{center}
}


\hypersetup{colorlinks=true,linkcolor=blue,citecolor=blue,urlcolor=blue}
\setlength{\emergencystretch}{3em}
\setlength{\tabcolsep}{5pt}
\renewcommand{\arraystretch}{1.08}

\newtheorem{proposition}{Proposition}[section]
\newtheorem{definition}[proposition]{Definition}
\newtheorem{remark}[proposition]{Remark}

\title{Toward a Particle-Spectrum Derivation from Observer-Overlap Consistency\\[0.4em]
\large A Status and Dependency Note for the OPH Particle Program}
\author{B.\ M\"uller}
\date{\OPHPaperReleaseDate}

\begin{document}
\maketitle
\OPHPaperReleaseBanner

\begin{abstract}
This note starts a dedicated public mathematical surface for the OPH particle
program. Its role is deliberately narrower than the synchronized compact paper:
it records the dependency boundary, separates calibration outputs from
theorem-level closure claims, and names the explicit missing objects in the
charged-lepton, quark, neutrino, and hadron lanes. The public boundary is as
follows. The D10 electroweak branch is a calibration-sector consistency
surface tied to the pixel-area input
\[
P \equiv a_{\mathrm{cell}}/\ell_P^2 = 1.63094,
\]
and exact closure for that lane is still open. The D11 Higgs/top critical-surface branch is
a secondary quantitative continuation once D10 is fixed. Charged leptons,
quarks, and neutrinos have explicit forward artifact boundaries and local guard
tests in the repository; they remain continuation lanes in the public ledger.
The hadron lane remains simulation-dependent
debug and systematics scaffolding only. The purpose of the note is therefore to
state the public rule cleanly: OPH exposes a structured particle
derivation program with sharpened constructive objects. It does not claim
first-principles closure of the full particle zoo. Accordingly, several public
mass rows still disagree with experiment, in some sectors substantially.
These mismatches mark open constructive work and should not be read as closed
derivations.
\end{abstract}

\section{Scope}

This note is a dependency and status surface for the OPH particle program. It
leaves the compact submission paper as the authoritative recovered-core
statement. A forward artifact in the code tree does not by itself promote a
sector.

\begin{definition}[Particle-program claim tiers]
The public particle program uses the following four-way split.
\begin{enumerate}[leftmargin=2em]
\item \textbf{Calibration sector:} an input-dependent quantitative branch used to
fix or audit the shared electroweak family once the allowed external input
\(P\) is declared.
\item \textbf{Secondary quantitative branch:} a branch that adds no new
continuous fit once the calibration sector is fixed, but still depends on
additional transport or matching structure outside the recovered core.
\item \textbf{Continuation lane:} a forward constructive lane with explicit
objects, artifact schemas, and guards, but without theorem-level closure.
\item \textbf{Simulation-dependent lane:} a lane whose outputs currently depend
on debug, toy, or systematics-oriented numerical scaffolding and therefore
cannot be promoted as a public OPH derivation claim.
\end{enumerate}
\end{definition}

\begin{proposition}[Current public particle boundary]
On the public repository snapshot represented by this note, only the D10
electroweak surface is
presented as a calibration sector and only the D11 Higgs/top surface is
presented as a secondary quantitative branch. The charged-lepton, quark,
neutrino, and hadron lanes remain below theorem-level closure.
\end{proposition}

\paragraph{Reason.}
The compact paper and the public repository ledger separate these lanes by
construction: the mirrored code tree carries explicit artifact boundaries and
guards for the open sectors, while the status ledger tags their rows as
continuation or simulation-dependent instead of structural or recovered-core
outputs.

\section{Dependency Ledger}

The particle program is best read as a staged map from the OPH core and
the declared external input \(P\) into increasingly ambitious quantitative
lanes. The table below is the public boundary that README, book, and downstream
publication surfaces should mirror.

\begingroup
\footnotesize
\setlength{\tabcolsep}{3pt}
\begin{longtable}{@{}>{\raggedright\arraybackslash}p{0.13\textwidth}>{\raggedright\arraybackslash}p{0.17\textwidth}>{\raggedright\arraybackslash}p{0.21\textwidth}>{\raggedright\arraybackslash}p{0.25\textwidth}>{\raggedright\arraybackslash}p{0.14\textwidth}@{}}
\caption{Public dependency ledger for the OPH particle program.}\\
\toprule
Lane & Public object & Immediate inputs & Smallest stated missing object & Public tier \\
\midrule
\endhead
\bottomrule
\endfoot
D10 electroweak & running-family electroweak closure & compact recovered core, declared \(P\), and printed running/matching conventions & coherent D10 electroweak transport kernel with shared scalar provenance across the correction family & calibration sector \\
D11 Higgs/top & synchronized critical-surface readout & D10 and the critical-surface synchronization branch & common D11 critical-surface readout kernel carrying the shared low-scale Higgs/top vector & secondary quantitative branch \\
Charged leptons & \(C_e^{\wedge} \to g_e \to Y_e\) forward chain & shared charged flavor dictionary and charged-budget transport artifact & theorem-level shared charged absolute-scale closure removing the remaining norm/gluing gap & continuation lane \\
Quarks/flavor & transport, projector, odd-response, and forward Yukawa artifacts & shared charged flavor lane, branch generator, and overlap-edge transport objects & theorem-level family excitation evaluator and closed odd quark response law removing the remaining constructive ambiguity & continuation lane \\
Neutrinos & capacity anchor, family response, Majorana lift, pullback metric, splittings, and PMNS gate & D6-style capacity anchor and shared flavor-basis artifacts & OPH-only scalar evaluator and downstream defect-weighted family needed for honest splitting/mixing promotion & continuation lane \\
Hadrons & quenched wrapper and systematics bookkeeping & legacy lattice reference script and explicit debug/systematics parameters & non-debug, non-quenched hadron closure with publishable systematics instead of smoke/demo artifacts & simulation-dependent lane \\
\end{longtable}
\endgroup

\section{Public Code Surface}

The repository carries a compact public mirror of the active particle code
under \texttt{code/particles/}. It exposes the current mathematical lanes
together with frozen public artifacts and tests. The full operational sandbox
stays outside this note.

The kept surfaces are:
\begin{enumerate}[leftmargin=2em]
\item the \texttt{core/} compatibility backbone used by the existing public
D10/D11 predictor surface;
\item the current lane directories
\texttt{calibration/}, \texttt{flavor/}, \texttt{leptons/},
\texttt{neutrino/}, and \texttt{hadron/};
\item the frozen artifact snapshot under \texttt{code/particles/runs/};
\item the public status ledger files \texttt{RESULTS\_STATUS.md},
\texttt{results\_status.json}, and \texttt{ledger.yaml} under
\texttt{code/particles/}.
\end{enumerate}

The intentionally omitted surfaces are Oracle batch archives, browser profiles,
autonomous supervisor plumbing, large run history, and the rest of the
non-public sandbox scaffolding. Active iteration continues in the sibling
\texttt{/particles} repo; this mirror exists so the public papers and summary
surfaces can link to a compact code snapshot.

\section{Mathematical Closure Program}

The particle program is organized around explicit constructive objects with
distinct claim tiers.

\subsection{D10 exactness}

The electroweak calibration lane still turns on transport/readout closure. The
public D10 artifact isolates the common-family
running electroweak surface, treats the active quintet as one shared
running-family object, and names the remaining burden as one
\(\mathrm{EWTransportKernel}_{D10}\) with a coherent shared scalar package
\(\Sigma_{EW,D10}\) beneath the pole/effective boundary. The open question is
transport/readout closure. More decimal places for \(P\) do not address that
issue.

\subsection{Charged sector}

The charged-lepton and quark lanes meet on a shared charged-budget surface. The
lepton side exposes the explicit
\texttt{charged\_common\_refinement\_mean\_eigenvalue\_certificate} and its
scalar-part subclause
\texttt{common\_refinement\_preserves\_mean\_eigenvalue}; the quark side
reduces the missing family-excitation law to a trace-zero quadratic evaluator
on the ordered branch-generator coordinate \(X_{\mathrm{ord}}\). This blocks
hidden split-scale fitting between leptons and quarks.

\subsection{Neutrinos}

The neutrino lane exposes the capacity anchor, the family-response tensor, the
Majorana lift, the pullback metric, the forward splitting artifacts, and the
defect-weighted \(\mu_e\) family as separate public objects. The sharp reduced
family beneath the open selector/evaluator closure is a positive raw edge score
\(q_e\) followed by the centered-log, mean-preserving lift to \(\mu_e\). The
remaining blocker is named sharply enough to fail.

\subsection{Hadrons}

The hadron branch remains the one lane where the repository is openly
simulation-dependent. The public wrapper makes that dependence explicit and
stores systematics/debug artifacts. Those artifacts stay in the
simulation-dependent tier and do not support a hadron-closure claim.

\section{Status Surfaces and Release Discipline}

The public repository exposes three synchronized particle-status surfaces:
\begin{enumerate}[leftmargin=2em]
\item the compact public code mirror in \texttt{code/particles/};
\item the frozen status ledger files \texttt{RESULTS\_STATUS.md} and
\texttt{results\_status.json} under \texttt{code/particles/};
\item the SVG overview used by the README and public summary pages.
\end{enumerate}

These surfaces are meant to stay synchronized with the compact paper and the
main README. If a lane is still tagged \emph{continuation} or
\emph{simulation-dependent} here, downstream book and web summaries should not
silently promote it.

\paragraph{Working rule.}
The public OPH particle story is a sharpened constructive program. It does not
claim a completed theorem-level derivation of the full particle spectrum. When
a public mass row is still numerically wrong, the named closure object for that
lane is still open. The published number has not earned promotion.

\end{document}
